% GENERATED FILE. Edit gujarati/pronunciation-reference.md, then run npm run generate:books.

\chapter*{The Gujarati script --- a reference}
\addcontentsline{toc}{chapter}{The Gujarati script (reference)}
\markboth{Gujarati script}{Gujarati script}

A \textbf{reference}, not a chapter — the lessons teach reading through real words. This page gathers the system and flags what is special about the \textbf{Gujarati script} and Gujarati sounds.

\section*{How the Gujarati script works (the three facts)}

\begin{enumerate}
\raggedright
\setlength{\itemsep}{0.35em}
  \item \textbf{Consonants carry a built-in \textquotedblleft{}a\textquotedblright{} — and there is no top line.} \gu{ક} = ka, \gu{ન} = na, \gu{મ} = ma, \gu{સ} = sa — written left to right. Gujarati is Devanagari's sister, but it \textbf{dropped the \emph{shirorekhā}} (the horizontal bar Devanagari, Bengali, and Gurmukhi hang their letters from) — cut it away, so its letters stand free. This missing top line is the one thing to remember, and the fastest way to tell Gujarati from Hindi at a glance — the \textquotedblleft{}headless script.\textquotedblright{}
  \item \textbf{A vowel sign changes the built-in \textquotedblleft{}a\textquotedblright{}:} \gu{કા} kā · \gu{કિ} ki · \gu{કી} kī · \gu{કુ} ku · \gu{કે} ke · \gu{કો} ko. The \emph{i}-sign \textbf{\gu{િ}} is written \emph{before} the consonant but read \emph{after} it. An \emph{anusvāra} \gu{ં} nasalises (\gu{છું} \emph{chhũ}); a \emph{chandra} \gu{ઁ} marks the light nasal on some words.
  \item \textbf{A \emph{virama} (\gu{્}) removes the vowel}, and the bare consonant joins the next as a conjunct: \gu{સ} + \gu{્} + \gu{ત} $\to$ \gu{સ્ત}.
\end{enumerate}

\section*{Standalone vowels (word-initial)}

Standalone vowels have their own letters, used when a word starts with a vowel: \gu{અ} a · \gu{આ} ā · \gu{ઇ} i · \gu{ઉ} u · \gu{એ} e · \gu{ઓ} o (\gu{આભાર} \emph{ābhār}).

\section*{What is special about Gujarati sounds and grammar}

\begin{itemize}
\raggedright
\setlength{\itemsep}{0.35em}
  \item \textbf{Three genders.} Gujarati keeps \textbf{Sanskrit's} masculine, feminine, \emph{and} neuter, visible on adjectives and possessives as \emph{-o / -ī / -ũ}: \emph{sāro} (m.) / \emph{sārī} (f.) / \emph{sārũ} (n.), \textquotedblleft{}good\textquotedblright{}; \emph{māro / mārī / mārũ}, \textquotedblleft{}my\textquotedblright{}. Hindi collapsed to two; Gujarati, like Marathi, kept three.
  \item \textbf{A copula all its own — \emph{chhe}.} Gujarati's \textquotedblleft{}is\textquotedblright{} is neither Hindi's \emph{hai} nor Sanskrit's \emph{asti}; it grew from an Old Gujarati form (\emph{achhaï}): \emph{hũ chhũ} (I am), \emph{tũ chhe}, \emph{te chhe}, \emph{tame chho}.
  \item \textbf{The retroflex \gu{ળ} (ḷ) — and \gu{ણ} (ṇ).} Gujarati keeps the extra \textquotedblleft{}hard l\textquotedblright{} — a \textbf{retroflex l}, tongue curled back — like Marathi and the Dravidian south, and absent from Hindi, as in \emph{maḷvũ} (\textquotedblleft{}to meet\textquotedblright{}). \textbf{\gu{ણ}} is the matching retroflex \emph{n}: the tongue tip curls back onto the roof of the mouth, as in \emph{jāṇvũ} (\textquotedblleft{}to know\textquotedblright{}). Middle Indo-Aryan regularly pulled a single \emph{n} between vowels back to this sound.
  \item \textbf{Aspirated stops are separate letters.} \textbf{\gu{ક}} \emph{ka} and \textbf{\gu{ખ}} \emph{kha} differ only by a puff of breath after the consonant — a difference English makes without noticing and Gujarati writes down (\emph{khāvũ}, \textquotedblleft{}to eat\textquotedblright{}).
  \item \textbf{Every verb is named in the neuter.} The infinitive ending \textbf{-\gu{વું}} \emph{-vũ} is the neuter of the three genders, so \emph{hovũ}, \emph{javũ}, \emph{khāvũ} all end the same way. Strip \emph{-vũ} and what is left is the stem.
  \item \textbf{Verb goes last; \textquotedblleft{}in\textquotedblright{} comes after the noun.} Gujarati is subject–object–verb and uses \textbf{post}positions: \emph{Amdāvad-mā rahũ chhũ} (\textquotedblleft{}I live in Ahmedabad\textquotedblright{}).
  \item \textbf{The trade-language layer.} Centuries as a great trading language across the Arabian Sea wove a heavy \textbf{Perso-Arabic} layer, and a few Portuguese words from the coastal ports, into everyday Gujarati — even \textquotedblleft{}how are you\textquotedblright{} is answered with the Persian loan \emph{majā} (\textquotedblleft{}enjoyment\textquotedblright{} $\leftarrow$ \emph{maza}).
\end{itemize}

\section*{Gujarati in the family}

Gujarati (\emph{gujarātī}) is \textbf{Indo-Aryan} — from Sanskrit, like Hindi, Marathi, and Punjabi — and its name comes from the \textbf{Gurjar}, a people who settled western India and gave their name to the land of \emph{Gujarāt}. Its Sanskrit core is shared with Hindi. Spoken by some 55 million people in Gujarat and across a vast trading diaspora, it is the mother tongue of \textbf{Mohandas K. (\textquotedblleft{}Mahatma\textquotedblright{}) Gandhi}. The lessons trace both threads — its Sanskrit roots and its Perso-Arabic trade layer — and, where they reach, the wider Indo-European family that also shaped English.
